\subsection{1. Introduction}\label{intro}

Fuel is the largest controllable operating cost of a commercial vehicle
fleet and a direct source of tailpipe emissions. When a vehicle
consistently burns more than expected, an operator faces a decision with
two very different remedies: coach the driver, or service the vehicle.
The remedies are not interchangeable. Sending an efficient driver to
remedial training wastes time and erodes trust, and dispatching a
mechanically sound vehicle to the workshop wastes a service slot;
conversely, a genuine fault left uncorrected keeps burning fuel and may
worsen. The decision therefore rests on knowing the cause of the excess,
not merely its size.

We treat this as an attribution problem. For a trip that consumed more
fuel than its route and conditions warrant, let the excess be the
difference between observed and expected consumption; the task is to
decompose that excess into a driver-behaviour component and a vehicle
component (the vehicle-baseline defined in Section 4.4). Raw fuel per
distance cannot perform this decomposition, because both an aggressive
driving style and a degraded engine inflate the same scalar. The
question is operationally valuable and, on public data, largely
unanswered: fuel-consumption models predict the scalar accurately
without decomposing its causes; eco-driving and behaviour-scoring work
quantifies the driver in isolation; and predictive-maintenance work
targets the fault directly but rarely treats fuel efficiency as the
symptom to explain.

Our approach rests on a physical observation about how the two causes
express themselves. Driver-caused excess is transient and kinematic: it
lives in acceleration, jerk, and harsh acceleration or braking events,
while the engine continues to run at its normal, near-stoichiometric
air-fuel ratio. Malfunction-caused excess, taking the rich-mixture fault
as the canonical fuel-relevant case, is steady-state and
combustion-linked: it shifts the air-fuel ratio away from stoichiometric
and leaves its trace in the exhaust chemistry, in elevated carbon
monoxide and unburnt hydrocarbons. Crucially, both signatures are
observable on standard on-board-diagnostics hardware. The kinematic axis
derives from GPS and inertial channels; the combustion axis derives from
the short- and long-term fuel trims that every closed-loop engine
reports, and, on a bench, from the air-fuel ratio directly.

The obstacle to demonstrating this on real data is that no public
dataset carries fuel consumption, mechanical-fault labels, and driver
behaviour together. A dedicated search (Section 3.4) confirms the gap:
the datasets that pair fuel with fault labels lack driver context or
ship only dimensionality-reduced features, and the driving-behaviour
datasets lack fuel. We therefore adopt a two-arm design on the datasets
that each carry one ground truth, and bridge them with a semi-synthetic
experiment that supplies the competing-cause labels no real dataset
provides. The fleet-telemetry arm (VED) grounds the decomposition and
the driver side; the engine-fault arm (EngineFaultDB) grounds the
malfunction signature; an independent labelled driving dataset validates
the kinematic axis; and the semi-synthetic injection measures
attribution accuracy directly.

The evidence chain (Figure 1) runs as follows: an expected-fuel model
and its variance decomposition establish that behaviour and a persistent
vehicle-baseline are separable, measurable shares of trip fuel; the
behaviour axis is validated to transfer to an independent dataset; the
rich-mixture fault is shown to produce measurable excess fuel with an
emissions signature; a signature contrast shows the combustion axis is
silent for driver-caused excess; and the semi-synthetic experiment shows
that, when fuel magnitude is made uninformative, signatures frozen from
their source domains recover the cause at 88.6\% accuracy (95.9\% for a
classifier trained on the injected labels), while fuel magnitude alone
is at chance. The remainder of the paper is organized as related work
(Section 2), data (Section 3), method and equations (Section 4),
experimental setup (Section 5), results for the two arms and the
attribution (Sections 6-8), discussion (Section 9), limitations (Section
10), and conclusion (Section 11).

\subsection{2. Related work}\label{related}

Four largely separate literatures meet in this problem: fuel-consumption
modelling, eco-driving and behaviour scoring, predictive maintenance and
engine-fault diagnosis, and telemetry anomaly detection. A fifth thread,
explainable attribution, contains the nearest prior work. We review each
and state where our contribution sits; the positioning is summarized in
Table 1.

\subsubsection{2.1 Fuel-consumption
modelling}\label{fuel-consumption-modelling}

Physics-based models map instantaneous kinematics to fuel and emissions.
Vehicle specific power {[}\hyperref[r7]{7}{]} expresses the tractive
power demand per unit mass and underlies many modal emission models; the
VT-Micro family {[}\hyperref[r8]{8}{]} fits polynomials in speed and
acceleration to instantaneous fuel and emission rates; and the
power-based VT-CPFM {[}\hyperref[r11]{11}{]} derives consumption from a
physical power formulation with few calibrated parameters. Data-driven
models now lead on on-board-diagnostics and telematics inputs: tree
ensembles are the reliable workhorse and have been shown to beat both
neural baselines and the engine-control-unit estimate against
ground-truth fuel meters {[}\hyperref[r4]{4}{]}, and kernel methods map
engine parameters to fuel rate from OBD-II {[}\hyperref[r9]{9}{]}. These
models predict consumption accurately, but they estimate a single
quantity; none decomposes that quantity into the driver and vehicle
contributions that an operator must separate.

\subsubsection{2.2 Eco-driving and driving-behaviour
scoring}\label{eco-driving-and-driving-behaviour-scoring}

Driving style moves fuel use substantially; data-driven eco-driving
reviews report savings of roughly 5-15\% from smoother driving
{[}\hyperref[r5]{5}{]}, a range our behaviour counterfactual reproduces
(Section 6.2), and recent reviews survey behaviour-to-fuel machine
learning specifically {[}\hyperref[r15]{15}{]}. Driver behaviour has
been classified from OBD signals and used to predict fuel consumption
{[}\hyperref[r12]{12}{]}, recognized as a driving style from smartphone
inertial sensors {[}\hyperref[r14]{14}{]}, and linked to efficiency
through drive-cycle simulation {[}\hyperref[r13]{13}{]}. This literature
scores the driver in isolation. It does not separate the
driver\textquotesingle s contribution from the
vehicle\textquotesingle s, and in single-driver-per-vehicle data the two
are confounded unless a model controls for one.

\subsubsection{2.3 Predictive maintenance and engine-fault
diagnosis}\label{predictive-maintenance-and-engine-fault-diagnosis}

Fault-diagnosis work detects sensor and actuator faults in the air-fuel
path from telemetry: neural detection and isolation of multiple
engine-sensor faults {[}\hyperref[r16]{16}{]}, fault-tolerant handling
of the mass-air-flow sensor whose signal we use to derive fuel
{[}\hyperref[r17]{17}{]}, and deep-learning fault diagnosis from driving
signals {[}\hyperref[r18]{18}{]}. The EngineFaultDB dataset we use as
our malfunction arm was released with a fault-classification baseline
{[}\hyperref[r2]{2}{]}; we repurpose it to quantify the fuel signature
of a fault rather than to classify faults. These methods target the
fault directly and treat fuel efficiency as an input, not as the symptom
to be attributed.

\subsubsection{2.4 Telemetry anomaly
detection}\label{telemetry-anomaly-detection}

The detector toolbox is mature: isolation forest
{[}\hyperref[r19]{19}{]}, one-class methods, and autoencoder or
recurrent models. In the automotive setting, however, most of this work
targets controller-area-network intrusion and safety rather than fuel
efficiency, for example LSTM autoencoders {[}\hyperref[r20]{20}{]},
recurrent autoencoders for embedded intrusion detection
{[}\hyperref[r21]{21}{]}, and spatiotemporal detection for connected
vehicles {[}\hyperref[r22]{22}{]}. These methods flag that a trip is
anomalous; they do not assign a fuel anomaly to a driver-versus -vehicle
cause. Closest in spirit, recent work detects anomalies in OBD snapshots
and adds a SHAP-based cause investigation with synthetically generated
anomalies {[}\hyperref[r24]{24}{]}; it investigates which features drive
a flagged anomaly, whereas we attribute the competing driver-versus
-vehicle cause across heterogeneous datasets that carry no joint
naturalistic cause labels. Fuel theft or leakage is a related but
distinct anomaly, detected as an abrupt change point
{[}\hyperref[r10]{10}{]}; that is a sudden step in fuel level, whereas
the excess we attribute is a sustained efficiency shift.

\subsubsection{2.5 Explainable attribution and the nearest prior
work}\label{explainable-attribution-and-the-nearest-prior-work}

Feature-attribution methods such as SHAP {[}\hyperref[r23]{23}{]}
explain a model\textquotesingle s output in terms of its input features.
The nearest prior work applies explainability to our exact problem:
Barbado and Corcho detect fuel-consumption anomalies on fleet telematics
with interpretable feature relevance and counterfactual recommendations
{[}\hyperref[r3]{3}{]}. Their method identifies which features co-occur
with an anomaly and what change would return an outlier to normal. It
attributes an anomaly to \emph{features}, not to a \emph{cause class}:
it does not say whether the excess is the driver or the vehicle, and its
feature relevances are not validated against an external ground truth
for either cause. Our signature-based decomposition is the increment: we
group observables into two physically motivated axes, validate each
against an independent signal, and measure the resulting cause
attribution. Methodologically the setting is a form of weak supervision
{[}\hyperref[r25]{25}{]}: the target fleet carries the outcome but no
cause labels, while each competing cause is labelled only in a separate
source domain with its own feature space, so cause-specific signatures
must be learned from the sources and combined on the target without
joint supervision.

\begin{longtable}[]{@{}
  >{\raggedright\arraybackslash}p{(\columnwidth - 8\tabcolsep) * \real{0.2000}}
  >{\raggedright\arraybackslash}p{(\columnwidth - 8\tabcolsep) * \real{0.2000}}
  >{\raggedright\arraybackslash}p{(\columnwidth - 8\tabcolsep) * \real{0.2000}}
  >{\raggedright\arraybackslash}p{(\columnwidth - 8\tabcolsep) * \real{0.2000}}
  >{\raggedright\arraybackslash}p{(\columnwidth - 8\tabcolsep) * \real{0.2000}}@{}}
\toprule\noalign{}
\begin{minipage}[b]{\linewidth}\raggedright
Approach
\end{minipage} & \begin{minipage}[b]{\linewidth}\raggedright
Detects excess
\end{minipage} & \begin{minipage}[b]{\linewidth}\raggedright
Ranks features
\end{minipage} & \begin{minipage}[b]{\linewidth}\raggedright
Attributes\\
driver vs vehicle\strut
\end{minipage} & \begin{minipage}[b]{\linewidth}\raggedright
Externally\\
validated\strut
\end{minipage} \\
\midrule\noalign{}
\endhead
\bottomrule\noalign{}
\endlastfoot
Fuel modelling + residual threshold {[}\hyperref[r4]{4}{]} & yes & no &
no & n/a \\
Eco-driving / behaviour scoring
{[}\hyperref[r5]{5},\hyperref[r12]{12}{]} & driver only & partial & no &
n/a \\
Engine-fault diagnosis {[}\hyperref[r16]{16}{]} & vehicle only & partial
& no & fault label \\
Anomaly detection + feature relevance {[}\hyperref[r3]{3}{]} & yes & yes
& no & no \\
\textbf{This work} & \textbf{yes} & \textbf{yes} & \textbf{yes
(measured, 88.6--95.9\%)} & \textbf{yes} \\
\end{longtable}

\textbf{Table 1.} Positioning against the four adjacent literatures and
the nearest prior work. Only the present method both attributes the
driver-versus-vehicle cause and validates its components externally.

\subsection{3. Data}\label{data}

\subsubsection{3.1 Vehicle Energy Dataset (driver
arm)}\label{vehicle-energy-dataset-driver-arm}

VED {[}\hyperref[r1]{1}{]} is a naturalistic driving dataset collected
in Ann Arbor, Michigan from November 2017 to November 2018, comprising
383 personal vehicles logged over roughly 374,000 miles through OBD-II
loggers, with synchronized GPS trajectories and engine channels. Each
record carries a timestamp, latitude and longitude, vehicle speed, mass
air-flow, engine RPM, absolute load, outside air temperature, auxiliary
power, and, for gasoline vehicles, the short- and long-term fuel trims;
records are keyed by a unique vehicle and trip identifier. We restrict
the analysis to the 264 internal-combustion (gasoline) vehicles,
excluding hybrids and electric vehicles whose fuel path is confounded by
the electric drive. The dedicated OBD fuel-rate parameter is essentially
unpopulated in VED, so fuel is derived from mass air-flow (Eq. 1), the
standard practice for this dataset. After resampling each trip to a
uniform 1 Hz grid and applying quality filters (a minimum trip distance
of 0.5 km, a minimum duration of 60 s, and at least 50\% valid mass
air-flow coverage), the pipeline yields \textbf{14,460 trips across 197
vehicles}. The variance decomposition and attribution analyses further
clip the fuel target at its 1st and 99th percentiles to remove sensor
glitches, leaving 14,170 trips across 194 vehicles; the semi-synthetic
study restricts to the 139 vehicles with at least 15 trips carrying
usable fuel-trim coverage, so that per-vehicle baselines are stable.
Median trip length is about 7 minutes.

\subsubsection{3.2 EngineFaultDB (malfunction
arm)}\label{enginefaultdb-malfunction-arm}

EngineFaultDB {[}\hyperref[r2]{2}{]} is a controlled test-bench dataset
built around a C14NE spark-ignition engine, with 55,999 records over 14
variables: manifold absolute pressure, throttle position, torque, power,
RPM, vehicle speed, exhaust concentrations of carbon monoxide,
hydrocarbons, carbon dioxide and oxygen, the air-fuel equivalence ratio
lambda, the air-fuel ratio, and fuel consumption in both litres per hour
and litres per 100 km. A fault label takes four values: 0 normal, 1 rich
mixture, 2 lean mixture, and 3 low-voltage ignition. The rich-mixture
label is confidently identified from its emission signature (elevated
carbon monoxide); the lean-versus-ignition mapping of labels 2 and 3 is
taken from the source paper\textquotesingle s listing order, as the
definitive mapping resides in a table published only as an image. This
uncertainty does not affect our result, which concerns the rich-mixture
fault. The bench carries no driver context, so it grounds the
malfunction fuel signature but not driving behaviour.

\subsubsection{3.3 Driving Events (behaviour
validation)}\label{driving-events-behaviour-validation}

The Driving Events dataset {[}\hyperref[r6]{6}{]} provides smartphone
inertial recordings from three drivers, with a linear-acceleration
stream at approximately 400 Hz and a companion file labelling time
windows as aggressive manoeuvres (right or left turn, braking,
acceleration, lane change) or non-aggressive events. After joining
sensor rows to labelled windows we obtain \textbf{169 events, 143
aggressive and 26 non-aggressive}. Two properties matter for
interpretation and are stated up front: the sensing regime differs from
VED (400 Hz smartphone inertial versus 1 Hz OBD kinematics), and the
class balance is skewed toward aggressive events.

\subsubsection{3.4 The joint-data gap}\label{the-joint-data-gap}

No public dataset carries fuel consumption, mechanical-fault labels, and
driver behaviour jointly, which is precisely what a single-dataset
attribution study would need. Our search of the candidates substantiates
this. EngineAD pairs a real-fleet fuel-rate channel with technician
fault labels but is released as eight principal components, so the raw
fuel signal is not recoverable, and access is gated. The diesel-fault
dataset DEFault carries fault labels but no fuel-consumption channel,
only in-cylinder pressure and vibration. SCANIA Component X provides
real heavy-truck maintenance records but fully anonymized signals with
no identifiable fuel variable, and the NASA C-MAPSS prognostics
benchmark is a simulated turbofan with no real malfunction and no driver
context. EngineFaultDB is the only openly downloadable dataset that
pairs fuel consumption with fault labels, and it is a bench without
drivers. This documented gap (Table 2) is the reason for the two-arm
design and the semi-synthetic bridge.

\begin{longtable}[]{@{}llllll@{}}
\toprule\noalign{}
Dataset & Role & Ground truth & Scale used & Sampling & License \\
\midrule\noalign{}
\endhead
\bottomrule\noalign{}
\endlastfoot
VED {[}\hyperref[r1]{1}{]} & Driver arm & Real fleet, behaviour & 197
cars / 14,460 trips & 1 Hz & Apache-2.0 \\
EngineFaultDB {[}\hyperref[r2]{2}{]} & Malfunction arm & Labelled engine
faults & 55,999 samples & bench & GPL-3.0 \\
Driving Events {[}\hyperref[r6]{6}{]} & Behaviour validation &
Aggressive vs normal & 3 drivers / 169 events & 400 Hz & CC-BY-4.0 \\
\end{longtable}

\textbf{Table 2.} The three public datasets, none of which carries fuel,
fault, and behaviour ground truth together.

\subsection{4. Method}\label{method}

\subsubsection{4.1 Problem formulation: attribution under missing joint
cause
labels}\label{problem-formulation-attribution-under-missing-joint-cause-labels}

We state the task abstractly, because the difficulty is general and not
specific to fuel. A target domain \(D_T\) contains an engineering
outcome \(y\) and telemetry \(x\) but \emph{no} labels for the competing
causes. Two competing causes, A (driver behaviour) and B (a vehicle
fault), are each labelled only in a separate source domain, \(D_A\) and
\(D_B\), which do not share the target\textquotesingle s raw feature
space. The goal is to attribute a target-domain excess to A or B, or to
abstain, without any jointly labelled example ever existing. Domain
knowledge supplies, for each cause, a transferable \emph{signature}: a
scoring function of the observables that rises when that cause is
present. The method (Algorithm 1) fits an expected-outcome model to
define excess on the target, learns each signature in its own source
domain and freezes it, places the two signatures on a common
standardized scale on the target, validates each against its source, and
attributes by comparing the two frozen scores with an abstention margin.
The remaining subsections instantiate each step for excess fuel; the
formulation itself is a form of weak supervision across heterogeneous
sources {[}\hyperref[r25]{25}{]}.

Algorithm 1 --- Multi-source signature attribution under missing joint
labels

\textbf{Input:} target \(D_T=\{(x_i,y_i)\}\) (no cause labels); sources
\(D_A,D_B\) labelling causes A, B; signature families \(s_A,s_B\);
abstention margin \(\tau\).\\
\textbf{Output:} per-unit cause \(\in\{A,B,\text{abstain}\}\) with a
confidence.

\begin{enumerate}

\item
  Fit expected-outcome model \(M\) on \(D_T\) (grouped OOF); excess
  \(e_i = y_i - M(x_i)\), grouped out-of-fold.
\item
  Learn signature \(s_A\) on \(D_A\) and \(s_B\) on \(D_B\);
  \textbf{freeze} both (no target labels used).
\item
  Score each target unit and standardize within the target:
  \(z^A_i, z^B_i\) (unsupervised).
\item
  Validate: confirm \(s_A\) separates cause A in \(D_A\) and \(s_B\)
  separates cause B in \(D_B\).
\item
  Attribute: \(A\) if \(z^A_i - z^B_i > \tau\); \(B\) if
  \(z^B_i - z^A_i > \tau\); else \textbf{abstain}. Confidence
  \(= |z^A_i - z^B_i|\).
\end{enumerate}

For excess fuel, \(M\) is the expected-fuel model (Section 4.3), cause A
is aggressive driving with a kinematic signature validated on labelled
driving (Section 6.3), and cause B is a rich-mixture fault with a
fueling/combustion signature validated on the engine bench (Section 7);
in the notation of Section 4.6, \(z^A = z^k\) and \(z^B = z^c\). The
frozen attribution and its abstention behaviour are evaluated in Section
8.4.

\subsubsection{4.2 Fuel target}\label{fuel-target}

Because the OBD fuel-rate parameter is unpopulated in VED, we derive the
instantaneous fuel rate from mass air-flow \(\dot m_{\text{air}}\) (g/s)
assuming a stoichiometric gasoline air-fuel ratio of 14.7 and a fuel
density of 745 g/L,

\begin{equation} \dot V_{\text{fuel}}\;[\text{L/h}] \;=\; \frac{\dot m_{\text{air}}}{14.7 \times 745}\times 3600, \end{equation}

and integrate over the 1 Hz trip grid, normalizing by trip distance to
obtain the trip target \(f\) in litres per 100 km. For the robustness
analysis of Section 8.5 we also form a trim-corrected target that
restores the fuel the stoichiometric assumption ignores,

\begin{equation} \dot V^{\text{corr}}_{\text{fuel}} \;=\; \dot V_{\text{fuel}}\left(1 + \frac{\text{STFT}+\text{LTFT}}{100}\right), \end{equation}

where STFT and LTFT are the short- and long-term fuel trims in percent.

\subsubsection{4.3 Expected-fuel model and
excess}\label{expected-fuel-model-and-excess}

We predict trip fuel from a behaviour feature block \(x^{\text{beh}}\)
and a conditions block \(x^{\text{env}}\) (Table 3) with a
gradient-boosted regression tree ensemble \(M\),

\begin{equation} \hat f_i \;=\; M\!\left(x^{\text{beh}}_i,\, x^{\text{env}}_i\right), \end{equation}

and define trip excess as the out-of-fold residual

\begin{equation} e_i \;=\; f_i - \hat f_i . \end{equation}

All predictions are out-of-fold under a GroupKFold split by vehicle, so
no vehicle is ever scored by a model that has seen any of its trips;
this prevents the model from memorizing a vehicle\textquotesingle s
identity and ensures the residual reflects deviation from the fleet, not
from the vehicle itself. Behaviour features include mean and dispersion
of speed, mean positive acceleration and its 95th percentile, jerk,
harsh-acceleration and harsh-braking rates (events exceeding
\(\pm 2.5\ \text{m/s}^2\), counted per km), idle fraction (speed below 3
km/h), stops per km, highway fraction, and vehicle specific power
computed with zero road grade; conditions include outside air
temperature, air-conditioning and heater power, vehicle weight,
distance, and duration.

\begin{longtable}[]{@{}ll@{}}
\toprule\noalign{}
Block & Features \\
\midrule\noalign{}
\endhead
\bottomrule\noalign{}
\endlastfoot
Behaviour \(x^{\text{beh}}\) & speed mean, speed SD, speed p85, mean
positive acceleration, acceleration p95, deceleration p05, jerk RMS,
harsh-accel/km, harsh-brake/km, idle fraction, stops/km, highway
fraction, vehicle specific power \\
Conditions \(x^{\text{env}}\) & outside air temperature, AC power,
heater power, weight, distance, duration \\
Combustion axis & short-term fuel trim, long-term fuel trim (VED);
air-fuel ratio, lambda, CO, HC (bench) \\
\end{longtable}

\textbf{Table 3.} Feature blocks. Harsh events use a
\(\pm 2.5\ \text{m/s}^2\) threshold; idle is speed below 3 km/h.
Engine-state channels (RPM, load) are deliberately excluded from the
expected-fuel model so that the residual is free to reflect vehicle
condition.

\subsubsection{4.4 Decomposition}\label{decomposition}

The persistent per-vehicle deviation, the \emph{vehicle-baseline}, is
the mean out-of-fold residual of vehicle \(j\),

\begin{equation} v_j \;=\; \frac{1}{n_j}\sum_{i \in j} e_i . \end{equation}

The \emph{behaviour component} of a trip is a counterfactual: the
model\textquotesingle s prediction for the observed driving minus its
prediction when the aggression features are set to a gentle reference
\(x^{\text{ref}}\) (the population 20th percentile of each aggression
feature), holding conditions fixed,

\begin{equation} b_i \;=\; M\!\left(x^{\text{beh}}_i, x^{\text{env}}_i\right) - M\!\left(x^{\text{ref}}, x^{\text{env}}_i\right). \end{equation}

To size the components at the population level we decompose the variance
of \(f\). With \(R^2_{\text{env}}\) the coefficient of determination of
a conditions-only model and \(R^2_{\text{full}}\) that of the full
model, the unique (marginal) behaviour share is
\(R^2_{\text{full}}-R^2_{\text{env}}\); the vehicle-baseline share is
the between-vehicle sum of squares of the residual expanded over trips,
divided by the total sum of squares,

\begin{equation} s_{\text{veh}} \;=\; \frac{\sum_i (v_{j(i)}-\bar e)^2}{\sum_i (f_i-\bar f)^2}, \end{equation}

and the unexplained share is \(1-R^2_{\text{full}}-s_{\text{veh}}\);
here \(\bar e\) and \(\bar f\) are the fleet-mean residual and
fleet-mean trip fuel.

\subsubsection{4.5 The two axes and the signature
contrast}\label{the-two-axes-and-the-signature-contrast}

The \emph{kinematic axis} summarizes transient aggression through the
behaviour features above; the \emph{combustion axis} measures the
deviation of the air-fuel ratio from stoichiometric. On the bench the
axis is the air-fuel-ratio deviation directly. In VED the observable is
the short- and long-term fuel trims, which are the engine control
unit\textquotesingle s feedback corrections to the base fuel schedule
rather than a direct air-fuel-ratio measurement; we therefore treat the
VED side as a fueling-correction signal,
\(c^{\text{VED}} = \text{STFT}+\text{LTFT}\), placed on a common
percent-fuel-per-air scale with the bench,

\begin{equation} c^{\text{bench}} \;=\; \left(\frac{14.7}{\text{AFR}}-1\right)\times 100 . \end{equation}

The signature contrast regresses each excess on each axis, reporting the
held-out coefficient of determination. The informative comparison is
whether the combustion axis, which by construction tracks fault-caused
excess, also tracks driver-caused excess; the claim is that it does not.

\subsubsection{4.6 Semi-synthetic
attribution}\label{semi-synthetic-attribution}

Because no real dataset labels a trip\textquotesingle s excess by cause,
we construct ground truth by injection. From each real base trip we form
a driver-caused and a fault-caused variant that raise fuel by the
\emph{same} amount \(\Delta f\), so that fuel magnitude carries no
information about the cause. The driver variant adds signal to the
kinematic axis and leaves the combustion axis at its real value; the
fault variant adds signal to the combustion axis and leaves kinematics
unchanged. The injected strength \(S\) is expressed in units of the
natural within-vehicle standard deviation of the axis, so that real
trip-to-trip variation supplies the confounding noise. Writing \(z^k\)
and \(z^c\) for the real within-vehicle-standardized kinematic and
combustion deviations,

\begin{equation} \text{driver:}\;(z^k+S,\;z^c), \qquad \text{fault:}\;(z^k,\;z^c+S). \end{equation}

We calibrate \(S\) to the measured rich-mixture fault: its combustion
shift of 3.0\% fuel-per-air, divided by the median within-vehicle
long-term-trim standard deviation of 1.79\%, gives
\(S_{\text{bench}} = 1.66\) (from the unrounded shift and standard
deviation). A logistic attributor is trained on the two axis-deviations
\((z^k, z^c)\) together with the vehicle baseline, under GroupKFold by
vehicle; a control classifier that sees only the fuel excess
\(\Delta f\) is the chance baseline.

\subsubsection{4.7 Baselines}\label{baselines}

We compare against the detection and feature-relevance methods the
literature offers: an isolation forest {[}\hyperref[r19]{19}{]} over
trip features as an unsupervised anomaly detector; permutation
importance over the fuel model as the feature-relevance surrogate used
by explainable anomaly work {[}\hyperref[r3]{3}{]}; and a pooled model
that omits any vehicle term, to show how much per-vehicle systematic
effect is left unattributed when the vehicle is not modelled.

\subsection{5. Experimental setup}\label{setup}

The expected-fuel model is a histogram gradient-boosted regression tree
ensemble (300 boosting iterations, learning rate 0.05, and maximum depth
5; the headline coefficient of determination is unchanged at 400
iterations and depth 6), trained under 5-fold GroupKFold by vehicle with
a fixed random seed. The fuel-model coefficient of determination is
stable across five seeds, ranging 0.672 to 0.676. Confidence intervals
throughout are cluster bootstraps that resample whole vehicles with
replacement over 1,000 draws and report the 2.5th and 97.5th
percentiles; for the behaviour-validation study the bootstrap resamples
events, and the permutation test uses 1,000 label permutations with the
leave-one-driver-out cross-validation repeated inside each. On the
bench, cross-validation must shuffle folds: the records are grouped by
fault and ordered within each group, so a non-shuffled 5-fold split
partitions the rich subset into disjoint operating ranges and the model,
forced to extrapolate, returns a spuriously negative coefficient of
determination; the rows are exchangeable for this regression, so
shuffled folds are correct. Software: Python 3.14, scikit-learn 1.8,
pandas 2.3, scipy 1.17, and pyarrow. Datasets are obtained from their
original sources and are not redistributed; the pipeline scripts
regenerate every number. The expected-fuel conclusions are not specific
to the gradient-boosted model: under the same vehicle-grouped protocol a
ridge baseline reaches \(R^2 = 0.663\) and a random forest 0.640,
against the boosted model\textquotesingle s 0.673 and a mean
predictor\textquotesingle s 0.00. Engine-state channels (RPM, load) are
deliberately excluded from the expected-fuel model; adding them raises
predictive \(R^2\) to 0.738 but they are engine outputs rather than
driver inputs, and their inclusion would let the predictor absorb part
of the diagnostic signature we aim to isolate.

\subsection{6. Results: driver arm (VED)}\label{arm1}

\subsubsection{6.1 Expected-fuel model}\label{expected-fuel-model}

The expected-fuel model reaches a held-out coefficient of determination
of \textbf{0.673} (95\% CI {[}0.628, 0.709{]}) with a mean absolute
percentage error of 13.3\% (95\% CI {[}12.3, 14.3{]}), stable across
seeds. This accuracy is in the range reported for OBD-based fuel models
on other fleets {[}\hyperref[r4]{4},\hyperref[r9]{9}{]}, though the
datasets differ and the numbers are not directly comparable; the
model\textquotesingle s role here is to define a well-calibrated
expectation from which excess is measured, not to set a new accuracy
record.

\subsubsection{6.2 Variance decomposition and the
vehicle-baseline}\label{variance-decomposition-and-the-vehicle-baseline}

Decomposing the variance of trip fuel (Table 4, Figure 2) shows that
route and environment dominate, that driving behaviour contributes a
unique 8\% beyond conditions and up to 37\% counting the variance it
shares with route, and that a persistent vehicle-baseline accounts for a
further 14\%. The vehicle-baseline is a real, size-independent effect:
regressing the per-vehicle fixed effect on engine displacement and
weight explains only 23\% of it, leaving roughly three-quarters
unexplained by vehicle size. This is the component that an operator
would want to route to inspection, and it is invisible to a model that
does not represent the vehicle.

The vehicle-baseline is not an in-sample artifact of estimating a group
mean from the same trips it is scored on. Estimating each
vehicle\textquotesingle s offset on one random half of its trips and
evaluating the fuel variance it explains on the held-out half gives a
cross-fitted share of \textbf{0.136--0.138} (stable across a minimum of
5, 10, 15, or 20 trips per vehicle), essentially unchanged from the
0.143 in-sample value, and an empirical-Bayes shrinkage estimate agrees
(0.138). The offsets are highly reliable: the split-half correlation of
per-vehicle baselines is \textbf{r = 0.95} (160 vehicles with at least
ten trips). A persistent, statistically stable per-vehicle component
therefore survives honest out-of-sample estimation. We label it a
persistent vehicle-associated residual, not a verified malfunction: VED
has effectively one driver per vehicle, so it cannot separate a
vehicle\textquotesingle s mechanical state from its habitual driver and
routes (Section 10); the malfunction side of the attribution is anchored
on the controlled bench, where the fault is actually induced.

\begin{longtable}[]{@{}lll@{}}
\toprule\noalign{}
Component & Share of fuel variance & 95\% CI \\
\midrule\noalign{}
\endhead
\bottomrule\noalign{}
\endlastfoot
Environment / route (conditions only, \(R^2_{\text{env}}\)) & 0.593 &
{[}0.542, 0.636{]} \\
Behaviour only (\(R^2_{\text{beh}}\)) & 0.368 & n/a \\
Full model (\(R^2_{\text{full}}\)) & 0.673 & {[}0.628, 0.709{]} \\
Behaviour, unique marginal & 0.080 & {[}0.060, 0.102{]} \\
Environment, unique marginal & 0.305 & n/a \\
Vehicle-baseline (\(s_{\text{veh}}\)) & 0.143 & {[}0.114, 0.174{]} \\
Unexplained & 0.184 & n/a \\
\end{longtable}

\textbf{Table 4.} Variance decomposition of trip fuel (14,170 trips, 194
vehicles), computed in one pass so the environment, unique-behaviour,
vehicle-baseline, and unexplained shares sum to one. Behaviour and route
share variance, so the unique 8\% understates
behaviour\textquotesingle s total association; the vehicle-baseline is a
separate 14\%. Point estimates; 95\% CIs are cluster bootstraps.

The behaviour counterfactual quantifies the driver component in physical
units. Relative to a gentle-driving reference, the median trip carries
an excess of \textbf{0.47 L/100km} (about 5\% of median fuel), rising to
1.59 at the 90th percentile. The component increases across deciles of
the harsh-event rate (Figure 3): its mean rises from 0.34 to 1.55
L/100km, monotone from the median upward, with a Spearman rank
correlation of 0.96 across deciles (\(p < 10^{-4}\)). Because the
harsh-event rate is itself an input to the behaviour features, this is
an internal face-validity check; the independent validation is the
external dataset of Section 6.3. The magnitude is consistent with the
5-15\% eco-driving range in the literature
{[}\hyperref[r5]{5},\hyperref[r15]{15}{]}.

\begin{figure}
\centering
\includegraphics[width=\linewidth]{fig1_variance_decomposition.png}
\caption{\textbf{Figure 2.} Variance decomposition of trip fuel. Route
and environment dominate; driving behaviour adds a unique 8\%, and a
persistent vehicle-baseline accounts for 14\%.}
\end{figure}

\begin{figure}
\centering
\includegraphics[width=\linewidth]{fig3_behaviour_monotonicity.png}
\caption{\textbf{Figure 3.} The counterfactual behaviour component rises
with independently measured aggression (Spearman \(\rho = 0.96\)).}
\end{figure}

\subsubsection{6.3 Behaviour-axis external
validation}\label{behaviour-axis-external-validation}

On the Driving Events dataset the kinematic amplitude features separate
aggressive from non-aggressive events with the driver held out.
Individually, the strongest features already discriminate
(acceleration-magnitude dispersion AUC 0.672, 95th percentile 0.657,
energy 0.646, harsh rate 0.639; all Mann-Whitney \(p<0.05\)), while raw
jerk at 400 Hz does not (AUC 0.48), consistent with amplitude rather
than derivative carrying the signal at that sampling rate. A logistic
classifier over the amplitude features reaches a leave-one-driver-out
AUC of \textbf{0.958} (95\% CI {[}0.927, 0.981{]}; per-driver 0.984,
0.934, 0.992; Figure 4). A random forest reaches 0.79; with only two
training drivers per fold the linear model generalizes across drivers
better, and we report both. Under a 1,000-permutation test the
label-permuted AUC falls to chance (mean 0.471), placing the real AUC at
\(p = 0.001\). This validates that the kinematic amplitude features
transfer to an independent accelerometer dataset; it is a different
sensing regime (400 Hz smartphone inertial versus 1 Hz OBD), so it
corroborates the axis rather than the exact VED feature set or the
counterfactual magnitude.

\begin{figure}
\centering
\includegraphics[width=\linewidth]{fig4_behaviour_roc.png}
\caption{\textbf{Figure 4.} Kinematic features separate aggressive from
normal driving on an independent dataset with the driver held out (AUC
0.958, 95\% CI {[}0.927, 0.981{]}).}
\end{figure}

To test transfer at the fleet\textquotesingle s own sampling rate and
across more drivers, we repeat the validation on a second, independent
behaviour dataset: the UAH-DriveSet {[}\hyperref[r26]{26}{]}, whose
processed streams are logged at \textbf{1 Hz} (matching VED) for
\textbf{six drivers} with normal, drowsy, and aggressive labels. Over
940 thirty-second windows (172 aggressive, 768 normal), kinematic
features computed at 1 Hz separate aggressive from normal driving with
the driver held out at \textbf{AUC 0.829} (logistic; 0.808 random
forest; per-driver 0.76--0.92; 1,000-permutation \(p = 0.001\)). The
lower value than the 400 Hz set is expected for the coarser sampling,
and it is the more deployment-relevant result: the behaviour axis
transfers to a different fleet, twice as many drivers, at the 1 Hz rate
OBD telematics actually provide.

\subsection{7. Results: malfunction arm (EngineFaultDB)}\label{arm2}

We fit the expected-fuel model on the normal (fault-0) operating points
only and score the faults at matched operating point. On normal data the
out-of-fold residual is centred at zero with a standard deviation of
0.20 L/100km, so the baseline is well calibrated. Scoring the faults
(Table 5, Figure 5), only the rich mixture burns measurably more:
\textbf{+0.50 L/100km, or +6.2\%} of median fuel, consistent with its
air-fuel ratio of 13.68 against the normal 14.07 (a 7.5\% enrichment of
fuel per unit air) and a brake-specific fuel consumption 10.6\% above
normal. The lean and low-voltage-ignition faults are not excess-fuel
faults. The rich-fault excess correlates with the combustion signature,
not with kinematics: with carbon monoxide 0.32, hydrocarbons 0.27, and
the air-fuel ratio \(-0.32\).

\begin{longtable}[]{@{}lllll@{}}
\toprule\noalign{}
Fault & Excess (L/100km) & \% of median fuel & AFR & BSFC vs normal \\
\midrule\noalign{}
\endhead
\bottomrule\noalign{}
\endlastfoot
Normal (0) & 0.00 & n/a & 14.07 & n/a \\
Rich mixture (1) & +0.50 & +6.2\% & 13.68 & +10.6\% \\
Lean mixture (2) & −0.06 & −0.7\% & 14.43 & n/a \\
Ignition (3) & +0.15 & +1.8\% & 14.39 & n/a \\
\end{longtable}

\textbf{Table 5.} Excess fuel at matched operating point by fault class.
Only the rich mixture is a fuel-relevant fault. Brake-specific fuel
consumption is reported for the rich fault only; the lean and ignition
faults were sampled at systematically lower power on the bench, so their
per-work figures are not comparable and are omitted.

Two checks guard the comparison. The source dataset reports that the
rich fault consumes less fuel \emph{per hour}; this reverses under our
matched analysis because the rich records were sampled at lower power,
so their per-hour consumption is low while their per-distance and
per-work consumption is high. Matching on operating point isolates the
true penalty. Second, the normal and rich samples overlap closely in
operating point (standardized mean differences: RPM 0.09, power 0.27,
speed 0.09), so the comparison is not driven by regime imbalance, and
the normal reference is itself moderately rich (lambda 0.957, carbon
monoxide 2.16\%), which biases the +6.2\% estimate conservatively
downward.

\begin{figure}
\centering
\includegraphics[width=\linewidth]{fig5_fault_signature.png}
\caption{\textbf{Figure 5.} Excess fuel at matched operating point by
fault class; only the rich mixture (low air-fuel ratio) burns measurably
more.}
\end{figure}

\subsection{8. Attribution}\label{attrib}

\subsubsection{8.1 Signature contrast}\label{signature-contrast}

Regressing each excess on each axis (Table 6, Figure 6) gives the
central qualitative result. The kinematic axis explains driver-caused
excess (held-out \(R^2 = 0.084\)) but is structurally unavailable on the
steady-state bench. The combustion axis explains fault-caused excess
strongly, but this cell is physically forced rather than discovered, and
the informative cell is its counterpart: the combustion axis is
\textbf{silent for driver-caused excess} (\(R^2 = -0.028\); univariate
correlation \(-0.106\)).

\begin{longtable}[]{@{}lll@{}}
\toprule\noalign{}
Axis & Driver-excess (VED) & Fault-excess (bench) \\
\midrule\noalign{}
\endhead
\bottomrule\noalign{}
\endlastfoot
Kinematic (transient aggression) & 0.084 & n/a (steady-state) \\
Combustion (AFR / fuel trim) & −0.028 & 0.903* \\
~~univariate correlation & −0.106 & +0.333 \\
\end{longtable}

\textbf{Table 6.} Signature contrast: variance of excess explained
(held-out \(R^2\)) by each axis. *The bench combustion cell is
physically forced (a rich mixture is extra fuel) and is not the claim;
see below.

Why the bench cell is not the claim. On the bench, air-fuel ratio and
fuel are physically coupled, so a combustion axis explaining
fault-excess is expected rather than discovered. The multivariate
\(R^2\) of 0.903 is carried mainly by carbon monoxide (\(R^2 = 0.73\))
and hydrocarbons (0.44), the unburnt-fuel emission products, with
air-fuel ratio and lambda perfectly collinear (\(r = 1.0\)) at 0.12
each; the univariate correlation of 0.33 understates the multivariate
fit because a single air-fuel-ratio transform misses the emission
channels. It is also interpolation, not generalization: under
leave-operating-region-out cross-validation (RPM×power terciles) the
all-sensor fit falls to \(R^2 = -0.48\) and a deployment-transferable
subset (air-fuel ratio, lambda, oxygen; excluding the bench-only CO and
HC analyzers) to 0.12, so the 0.903 reflects near-neighbour operating
points rather than a portable predictor. We therefore treat this cell as
a physical consistency check and rest the claim on the silent
driver-excess cell and the attribution accuracy below. Of the
\(2\times2\), one cell is physically forced, one is structurally absent,
and the two evidential cells are the silent combustion axis on
driver-excess and the kinematic axis carrying the driver signal; we
therefore call this a signature contrast rather than a full double
dissociation.

\begin{figure}
\centering
\includegraphics[width=\linewidth]{signature_contrast.png}
\caption{\textbf{Figure 6.} Signature contrast. The combustion axis is
silent for driver-caused excess while the kinematic axis carries it; the
starred bench cell is physically expected.}
\end{figure}

\subsubsection{8.2 Comparison to detection and feature-relevance
baselines}\label{comparison-to-detection-and-feature-relevance-baselines}

Existing methods detect and rank but do not attribute. The isolation
forest tracks excess only weakly (correlation 0.16 with excess), aligns
with the behaviour component (0.22) but not the vehicle-baseline
(\(-0.01\)), and returns a single anomaly score with no cause label.
Permutation-importance feature relevance ranks mean speed (0.95) and
weight (0.69) highest, physical determinants of consumption but not a
per-trip driver-versus-vehicle split. A pooled model that omits the
vehicle term leaves 1.23 L/100km of per-vehicle systematic effect
unattributed. None of these produces the decomposition our method does;
because there is no cause-labelled real benchmark, this is a capability
contrast, and the quantitative evaluation is next.

\subsubsection{8.3 Attribution accuracy on semi-synthetic ground
truth}\label{attribution-accuracy-on-semi-synthetic-ground-truth}

Injecting driver-caused and fault-caused excess into 11,520 real VED
trips (139 vehicles) with the same fuel increase for both, so that
magnitude is uninformative, the signature attributor recovers the cause
with accuracy that rises with signature strength: 0.713 at \(S=0.5\),
0.866 at 1.0, 0.946 at 1.5, 0.979 at 2.0, and 0.998 at 3.0, while the
fuel-only baseline stays at chance (about 0.50) throughout (Figure 7).
At the bench-calibrated magnitude \(S_{\text{bench}}\approx 1.66\) the
accuracy is \textbf{0.959} (95\% CI {[}0.956, 0.962{]}, F1 0.959, with
near-symmetric errors of about 4\% each way). This measures attribution
rather than asserting it: the design makes fuel magnitude uninformative
by construction, so the accuracy comes entirely from the signature, and
by controlling the injected magnitude it sidesteps the trim-blindness of
the fleet fuel target and the healthy-fleet range restriction examined
next. Two caveats on scope. The injection places each cause purely on
its own axis, so the experiment measures how detectable that separation
is against real trip-to-trip noise; it presumes rather than proves
axis-separability, which rests instead on the bench signature and the
silent driver-excess cell. And the measured accuracy sits above the
Gaussian-optimal bound of \(\Phi(S/\sqrt 2)\approx 0.88\) at
\(S_{\text{bench}}\), because the real within-vehicle deviations are
tighter than the median standard deviation used for calibration and the
attributor also sees each vehicle\textquotesingle s baseline; the two
axes are in fact near-independent in the fleet (trip-level correlation
\(-0.01\)), which is what makes them separable.

\begin{figure}
\centering
\includegraphics[width=\linewidth]{fig6_attribution_accuracy.png}
\caption{\textbf{Figure 7.} Attribution accuracy versus signature
strength, reaching 95.9\% at the bench-calibrated fault magnitude while
the fuel-magnitude baseline stays at chance. This classifier is trained
on the injected labels and is therefore an upper bound; Section 8.4
reports the transfer result that does not use them.}
\end{figure}

\subsubsection{8.4 Transfer with frozen source-derived
signatures}\label{transfer-with-frozen-source-derived-signatures}

The classifier above is trained on the same injection labels it
recovers, so it measures recoverability rather than transfer. We
therefore repeat the evaluation with \textbf{frozen} signatures that
never see the VED injection labels: the driver signature is a logistic
model learned only on the external aggressive-driving dataset (over
kinematic features matched to VED: acceleration variability, its 95th
percentile, and harsh-event rate), and the fault signature is the
combustion axis, whose direction and effect size come from the bench.
Each VED trip receives a frozen kinematic score and a frozen combustion
score, and the attribution rule is simply which standardized score is
larger, a fixed rule with no fitted parameters. At the bench-calibrated
magnitude this zero-shot rule reaches \textbf{88.6\% accuracy} (Table
7), matching a logistic classifier trained on the injected labels over
the same two scores (0.886 vs 0.886) and far above the fuel-only chance
level; the full-feature classifier of Section 8.3 sets the upper bound
at 95.9\%, and the frozen and same-score curves track across effect
sizes (0.86 at S = 1.5, 0.92 at S = 2.0). Because the scoring functions
are learned in their source domains and combined on the target without
joint supervision, this is the multi-source transfer the framework
claims, not a rediscovery of the injection rule.

The frozen rule also supports abstention. Withholding a decision when
the two scores are close (a margin \(\tau\)) trades coverage for
accuracy: at the bench magnitude, deciding on the 71\% of trips with the
clearest evidence raises accuracy from 88.6\% to 95.7\%, and on the
clearest 41\% to 98.2\%. In a maintenance setting an explicit
"uncertain" verdict is preferable to a confident wrong intervention.

\begin{longtable}[]{@{}llll@{}}
\toprule\noalign{}
Signature strength S & Frozen transfer & Trained on same scores &
Fuel-only \\
\midrule\noalign{}
\endhead
\bottomrule\noalign{}
\endlastfoot
1.0 & 0.769 & 0.768 & 0.50 \\
1.5 & 0.863 & 0.863 & 0.50 \\
1.66 (bench) & 0.886 & 0.886 & 0.50 \\
2.0 & 0.924 & 0.925 & 0.50 \\
3.0 & 0.980 & 0.981 & 0.51 \\
\end{longtable}

\textbf{Table 7.} Attribution accuracy by signature strength. The frozen
source-derived rule (no injection-label training) matches a classifier
trained on the same two scores with the injection labels; fuel magnitude
alone stays at chance.

The real world does not present pure one-axis changes, so we stress the
frozen rule. When both causes act at once, we score dominant-cause
accuracy over a grid of driver and fault effect sizes (Figure 8):
attribution is accurate off the diagonal and degrades to chance only
where the two effects are equal and the cause is genuinely ambiguous.
The rule degrades gracefully under departures from ideal conditions at
the bench magnitude: additive Gaussian score noise takes accuracy from
0.886 to 0.80 (1 SD) and 0.74 (1.5 SD); a systematic combustion-score
bias to 0.83 (1 SD); and dropping the fuel-trim channel for a fraction
of trips to 0.85 (20\% missing) and 0.80 (50\%). As a probabilistic
scorer (the sigmoid of the score margin) the frozen rule reaches AUROC
0.951, average precision 0.945, and a Brier score of 0.092 at the bench
magnitude.

\begin{figure}
\centering
\includegraphics[width=\linewidth]{mixed_cause_grid.png}
\caption{\textbf{Figure 8.} Dominant-cause accuracy when both causes
act, over driver effect \(S_d\) and fault effect \(S_f\). Attribution is
reliable off the diagonal and falls to chance only where the two effects
are equal.}
\end{figure}

Attribution is not an artifact of one classifier. At the bench
magnitude, task-aligned baselines all cluster near the frozen rule and
far above chance (Table 8): a raw-threshold rule on harsh-event rate
versus fuel trim reaches 0.89, a likelihood-ratio prototype against
source-calibrated signature distributions 0.89, and a classifier trained
on all features and the injection labels sets the upper bound at 0.95,
while the fuel-only baseline stays at 0.51. The signature separation,
not the choice of attributor, is what carries the result.

\begin{longtable}[]{@{}lll@{}}
\toprule\noalign{}
Attributor (at bench magnitude) & Accuracy & Uses injection labels? \\
\midrule\noalign{}
\endhead
\bottomrule\noalign{}
\endlastfoot
Fuel magnitude only & 0.51 & n/a (chance) \\
Frozen source-derived argmax & 0.886 & no \\
Raw-threshold rule (harsh rate vs trim) & 0.891 & no \\
Likelihood-ratio prototype (source-calibrated) & 0.886 & no \\
Trained full-feature logistic (upper bound) & 0.952 & yes \\
\end{longtable}

\textbf{Table 8.} Task-aligned attribution baselines. Every
signature-based attributor reaches \textasciitilde0.89 without injection
labels; only the trained upper bound uses them; fuel magnitude is at
chance.

\subsubsection{8.5 Robustness of the silent-combustion
result}\label{robustness-of-the-silent-combustion-result}

Two threats could make the silent combustion cell an artifact rather
than a finding; both are ruled out.

\paragraph{Trim-blindness of the fuel
target}\label{trim-blindness-of-the-fuel-target}

The fleet fuel target is derived from mass air-flow at stoichiometric
ratio (Eq. 1), which ignores the trim-driven fuel fraction and could
suppress the combustion axis by construction. Re-deriving fuel with the
trim correction of Eq. 2 moves the combustion axis on driver-excess only
from \(R^2 = -0.028\) to \(+0.012\); it stays silent, because the
ignored fraction is small in this healthy fleet (median 1.5\%, 90th
percentile 11\%). The result is also robust to how the
fueling-correction signal is treated: the short- and long-term trims are
individually uninformative for driver-excess (\(R^2 = 0.019\) and
\(-0.051\)), and an additive versus multiplicative trim correction
changes the fuel target by a median of 0.002 L/100km and the combustion
\(R^2\) by less than 0.01.

\paragraph{Range restriction}\label{range-restriction}

If VED simply had no combustion variance, the axis could not correlate
with anything. It does have variance (within-vehicle long-term-trim
standard deviation 1.79\%, between-vehicle 4.41\%). Restricting to the
highest-trim quartile of vehicles (47 vehicles) still gives a near-zero
excess-versus-trim correlation (\(r = -0.14\)). The silence therefore
reflects that driver-caused excess is genuinely not combustion-linked,
not a measurement or sampling artifact.

\subsubsection{8.6 A worked example}\label{a-worked-example}

Figure 9 shows the method\textquotesingle s per-trip product on two real
trips with the same symptom and opposite causes. Trip A burns 5.2
L/100km above its gentle-driving baseline, almost all of it driving
behaviour, with a negligible vehicle-baseline; Trip B burns 2.8 above
baseline, almost all of it a persistent vehicle-baseline offset, with
negligible behaviour. Fuel per distance alone would flag both and
explain neither; the decomposition attributes each to its cause.

\begin{figure}
\centering
\includegraphics[width=\linewidth]{fig7_worked_example.png}
\caption{\textbf{Figure 9.} Two real trips, same elevated fuel, opposite
cause. The decomposition separates a driver-dominated trip (A) from a
vehicle-dominated trip (B).}
\end{figure}

\subsection{9. Discussion}\label{discuss}

\subsubsection{9.1 Why signature-based attribution
works}\label{why-signature-based-attribution-works}

Fuel per distance is a sum over causes, and summation destroys the
information needed to separate them: two trips with identical excess can
arise from opposite mechanisms. The signature survives because each
cause perturbs the channel it physically acts on. Aggressive driving
raises the tractive power demand while the engine holds a normal
air-fuel ratio, so the excess appears in kinematics and not in the
trims; a rich mixture over-fuels at a given air-flow, so the excess
appears in the air-fuel ratio and the unburnt-fuel emissions and not in
kinematics. In a healthy fleet the two axes are nearly orthogonal
(trip-level correlation \(-0.01\)), which is exactly why the combustion
axis is silent for driver-caused excess and why attribution is possible
from standard channels.

\subsubsection{9.2 Deployment}\label{deployment}

Both axes come from signals a fleet telematics unit already collects:
GPS and inertial derivatives for the kinematic axis, and the short- and
long-term fuel trims for the combustion axis. The per-trip decomposition
can therefore be computed online and routed: a persistent
vehicle-baseline offset feeds an inspection queue, while a large
behaviour component feeds driver coaching, in both cases with the
attribution\textquotesingle s uncertainty attached and a human in the
loop. The accuracy curve of Section 8.3 sets expectations for smaller
faults: real degradations weaker than the bench rich fault correspond to
\(S<1.66\) and are attributed less confidently, degrading gracefully
rather than failing abruptly.

\subsubsection{9.3 What the vehicle-baseline
is}\label{what-the-vehicle-baseline-is}

The vehicle-baseline is a real 14\% of trip-fuel variance,
three-quarters of it unexplained by engine size or weight. Its mechanism
is not labelled in VED and plausibly bundles tuning, tire pressure,
accumulated wear, and accessory load. That VED cannot resolve this
further is itself informative: the fleet is healthy, which is why the
malfunction side of the attribution is anchored on a controlled bench
rather than observed in the fleet. Attribution here identifies
\emph{that} a vehicle deviates persistently, not the specific mechanical
fault, which remains a diagnosis step.

\subsubsection{9.4 The benchmark the field
needs}\label{the-benchmark-the-field-needs}

Naming what would close the joint-data gap (Section 3.4) turns it into
an agenda: a fleet logged with true fuel-flow meters, OBD trims, driver
identifiers, and verified maintenance or fault records on the same
vehicles and trips. Such a dataset would allow single-trip
competing-cause attribution to be trained and measured directly, and the
semi-synthetic bridge to be replaced by real ground truth.

\subsection{10. Scope and limitations}\label{limits}

Driver and fault ground truths live in different datasets, so the
signature contrast is cross-dataset; the semi-synthetic experiment
provides the single-trip competing-cause evaluation that no real dataset
supports, but it assumes additive, axis-separable effects, and a real
fault that co-moved both axes would make attribution harder than the
controlled estimate. VED has effectively one driver per vehicle, so
per-trip driver-versus-vehicle causation is not identifiable from
observation alone; the decomposition is model-based and its attribution
is validated on injected ground truth. The malfunction arm is a
controlled test bench with no driver context, and its fuel-relevant
fault is the rich mixture; the lean-versus-ignition label mapping
follows the source paper. VED is a single-city, single-climate,
gasoline-only fleet, and its telematics are resampled to 1 Hz, so fine
transient dynamics are smoothed; the behaviour-axis transfer is shown
for amplitude features across sensing regimes rather than for the exact
VED feature set, and a dedicated 1 Hz-versus-high-rate sensitivity
analysis is left to future work.

\subsubsection{10.1 Responsible use}\label{responsible-use}

Driver-side attribution is, in fleet practice, a form of worker
monitoring, and an attribution error that blames a driver for a
mechanical fault, or the reverse, has real consequences for pay and
employment. The method should be used with calibrated uncertainty and
human review, never as an automatic verdict; the per-trip output is a
decomposition with error bars, not a judgment.

\subsection{11. Conclusion}\label{conclusion}

Excess fuel consumption has two causes with one symptom, and a
fuel-per-distance figure cannot tell a heavy-footed driver from a
thirsty vehicle. We showed that the causes separate by signature: driver
excess is transient and kinematic, malfunction excess is steady-state
and combustion-linked. On a 194-vehicle public fleet, behaviour and a
persistent vehicle-baseline account for measurable 8\% and 14\% shares
of trip-fuel variance; on a controlled bench, a rich-mixture fault burns
6.2\% more with an emissions signature; the combustion axis is silent
for driver-caused excess, and that silence survives both a
trim-corrected fuel target and the highest-trim vehicles. When fuel
magnitude is made uninformative, frozen signatures learned in separate
source domains attribute the cause at 88.6\% in a calibrated
semi-synthetic transfer test (trained upper bound 95.9\%), where
magnitude alone is at chance. Because both axes derive from standard
on-board-diagnostics channels, per-trip attribution is deployable on
existing telematics, and the main obstacle to pushing it further is the
absence of a public dataset that carries fuel, faults, and behaviour
together.

\subsection{Data and code availability}\label{repro}

Datasets are obtained from their original sources and not redistributed:
VED (\href{https://github.com/gsoh/VED}{github.com/gsoh/VED},
Apache-2.0); EngineFaultDB
(\href{https://github.com/leoxthomas/EngineFaultDB}{github.com/leoxthomas/EngineFaultDB},
GPL-3.0); Driving Events
(\href{https://zenodo.org/records/6570972}{zenodo.org/records/6570972},
CC-BY-4.0). The full pipeline (\texttt{scripts/p1\_features.py} through
\texttt{scripts/p14\_stress\_baselines.py}), results log
(\texttt{RESULTS.md}), and this manuscript are in the project
repository; every number traces to a script. Environment: Python 3.14,
scikit-learn 1.8, pandas 2.3, scipy 1.17, pyarrow.

\textbf{Author contributions.} A. Apartsin conceived and implemented the
study and drafted the manuscript; Y. Aperstein contributed to the study
design and critically revised the manuscript. \textbf{Acknowledgements.}
We thank the creators of VED, EngineFaultDB, and the Driving Events
dataset for releasing their data openly.

\subsection{References}\label{refs}

\begin{enumerate}

\item
  \phantomsection\label{r1}{G. Oh, D. J. LeBlanc, H. Peng. Vehicle
  Energy Dataset (VED), A Large-scale Dataset for Vehicle Energy
  Consumption Research. \emph{IEEE Transactions on Intelligent
  Transportation Systems}, 2020.
  \href{https://doi.org/10.1109/TITS.2020.3035596}{10.1109/TITS.2020.3035596}}
\item
  \phantomsection\label{r2}{L. Thomas et~al. EngineFaultDB: A Novel
  Dataset for Automotive Engine Fault Classification and Baseline
  Results. \emph{IEEE Access}, 2023.
  \href{https://doi.org/10.1109/ACCESS.2023.3331316}{10.1109/ACCESS.2023.3331316}}
\item
  \phantomsection\label{r3}{A. Barbado, Ó. Corcho. Interpretable machine
  learning models for predicting and explaining vehicle fuel consumption
  anomalies. \emph{Engineering Applications of Artificial Intelligence},
  2022.
  \href{https://doi.org/10.1016/j.engappai.2022.105222}{10.1016/j.engappai.2022.105222}}
\item
  \phantomsection\label{r4}{H. Abediasl, A. Ansari, V. Hosseini, C. R.
  Koch, M. Shahbakhti. Real-time vehicular fuel consumption estimation
  using machine learning and on-board diagnostics data. \emph{Proc.
  IMechE Part D}, 2024.
  \href{https://doi.org/10.1177/09544070231185609}{10.1177/09544070231185609}}
\item
  \phantomsection\label{r5}{P. Fafoutellis, E. G. Mantouka, E. I.
  Vlahogianni. Eco-Driving and Its Impacts on Fuel Efficiency: An
  Overview of Technologies and Data-Driven Methods.
  \emph{Sustainability}, 13(1):226, 2021.
  \href{https://doi.org/10.3390/su13010226}{10.3390/su13010226}}
\item
  \phantomsection\label{r6}{J. Vieira et~al. Driving Events Dataset.
  Zenodo, 2022.
  \href{https://doi.org/10.5281/zenodo.6570972}{10.5281/zenodo.6570972}}
\item
  \phantomsection\label{r7}{J. L. Jimenez-Palacios. Understanding and
  quantifying motor vehicle emissions with vehicle specific power and
  TILDAS remote sensing. PhD thesis, MIT, 1999.}
\item
  \phantomsection\label{r8}{K. Ahn, H. Rakha, A. Trani, M. Van Aerde.
  Estimating Vehicle Fuel Consumption and Emissions based on
  Instantaneous Speed and Acceleration Levels. \emph{Journal of
  Transportation Engineering}, 128(2):182--190, 2002.
  \href{https://doi.org/10.1061/(ASCE)0733-947X(2002)128:2(182)}{10.1061/(ASCE)0733-947X(2002)128:2(182)}}
\item
  \phantomsection\label{r9}{T. Abukhalil, H. AlMahafzah, M. Alksasbeh,
  B. A. Y. Alqaralleh. Fuel Consumption Using OBD-II and Support Vector
  Machine Model. \emph{Journal of Robotics}, 2020:9450178.
  \href{https://doi.org/10.1155/2020/9450178}{10.1155/2020/9450178}}
\item
  \phantomsection\label{r10}{R. Chu, W. H. Chik, J. Song, J. Chan, X.
  Li. Real-time Fuel Leakage Detection via Online Change Point
  Detection. arXiv:2410.09741, 2024.
  \href{https://doi.org/10.48550/arXiv.2410.09741}{10.48550/arXiv.2410.09741}}
\item
  \phantomsection\label{r11}{H. A. Rakha, K. Ahn, K. Moran, B. Saerens,
  E. Van den Bulck. Virginia Tech Comprehensive Power-Based Fuel
  Consumption Model (VT-CPFM). \emph{Transportation Research Part D},
  16(7):492--503, 2011.
  \href{https://doi.org/10.1016/j.trd.2011.05.008}{10.1016/j.trd.2011.05.008}}
\item
  \phantomsection\label{r12}{P. Ping, W. Qin, Y. Xu, C. Miyajima, K.
  Takeda. Impact of Driver Behavior on Fuel Consumption: Classification,
  Evaluation and Prediction Using Machine Learning. \emph{IEEE Access},
  7, 2019.
  \href{https://doi.org/10.1109/ACCESS.2019.2920489}{10.1109/ACCESS.2019.2920489}}
\item
  \phantomsection\label{r13}{A. Mohammadnazar, Z. H. Khattak, A. J.
  Khattak. Assessing driving behavior influence on fuel efficiency using
  machine learning and drive-cycle simulations. \emph{Transportation
  Research Part D}, 126:104025, 2024.
  \href{https://doi.org/10.1016/j.trd.2023.104025}{10.1016/j.trd.2023.104025}}
\item
  \phantomsection\label{r14}{D. A. Johnson, M. M. Trivedi. Driving style
  recognition using a smartphone as a sensor platform. \emph{IEEE ITSC},
  2011.
  \href{https://doi.org/10.1109/ITSC.2011.6083078}{10.1109/ITSC.2011.6083078}}
\item
  \phantomsection\label{r15}{S. S. Haghshenas, A. Kouvelas, V. Astarita,
  S. S. Haghshenas, G. Guido. Driver behavior classification and its
  impact on fuel consumption using machine learning: A comprehensive
  review. \emph{Transportation Engineering}, 2026.
  \href{https://doi.org/10.1016/j.treng.2026.100436}{10.1016/j.treng.2026.100436}}
\item
  \phantomsection\label{r16}{M. Cervantes-Bobadilla et~al. Multiple
  fault detection and isolation using artificial neural networks in
  sensors of an internal combustion engine. \emph{Engineering
  Applications of AI}, 117:105524, 2023.
  \href{https://doi.org/10.1016/j.engappai.2022.105524}{10.1016/j.engappai.2022.105524}}
\item
  \phantomsection\label{r17}{P. Gutiérrez León et~al. Implementation of
  a fault tolerant system for the internal combustion
  engine\textquotesingle s MAF sensor. \emph{Measurement}, 122, 2018.
  \href{https://doi.org/10.1016/j.measurement.2018.03.006}{10.1016/j.measurement.2018.03.006}}
\item
  \phantomsection\label{r18}{C.-S. A. Gong, C.-H. S. Su, Y.-E. Liu,
  D.-Y. Guu, Y.-H. Chen. Deep Learning with LPC and Wavelet Algorithms
  for Driving Fault Diagnosis. \emph{Sensors}, 22(18):7072, 2022.
  \href{https://doi.org/10.3390/s22187072}{10.3390/s22187072}}
\item
  \phantomsection\label{r19}{F. T. Liu, K. M. Ting, Z.-H. Zhou.
  Isolation Forest. \emph{IEEE ICDM}, 2008.
  \href{https://doi.org/10.1109/ICDM.2008.17}{10.1109/ICDM.2008.17}}
\item
  \phantomsection\label{r20}{S. Longari, D. H. Nova Valcarcel, M. Zago,
  M. Carminati, S. Zanero. CANnolo: An Anomaly Detection System Based on
  LSTM Autoencoders for Controller Area Network. \emph{IEEE Trans.
  Network and Service Management}, 18(2), 2021.
  \href{https://doi.org/10.1109/TNSM.2020.3038991}{10.1109/TNSM.2020.3038991}}
\item
  \phantomsection\label{r21}{V. K. Kukkala, S. V. Thiruloga, S.
  Pasricha. INDRA: Intrusion Detection Using Recurrent Autoencoders in
  Automotive Embedded Systems. \emph{IEEE Trans. CAD of Integrated
  Circuits and Systems}, 39(11), 2020.
  \href{https://doi.org/10.1109/TCAD.2020.3012749}{10.1109/TCAD.2020.3012749}}
\item
  \phantomsection\label{r22}{P. Mansourian, N. Zhang, A. Jaekel, M.
  Kneppers. Deep Learning-Based Anomaly Detection for Connected
  Autonomous Vehicles Using Spatiotemporal Information. \emph{IEEE
  Trans. Intelligent Transportation Systems}, 2023.
  \href{https://doi.org/10.1109/TITS.2023.3286611}{10.1109/TITS.2023.3286611}}
\item
  \phantomsection\label{r23}{S. M. Lundberg, S.-I. Lee. A Unified
  Approach to Interpreting Model Predictions (SHAP). arXiv:1705.07874,
  2017.
  \href{https://doi.org/10.48550/arXiv.1705.07874}{10.48550/arXiv.1705.07874}}
\item
  \phantomsection\label{r24}{V. Vučinić et~al. Anomaly Detection for
  Vehicle Diagnostics Based on OBD Snapshots With Cause Investigation.
  \emph{IEEE Access}, 14:24765--24785, 2026.
  \href{https://doi.org/10.1109/ACCESS.2026.3664118}{10.1109/ACCESS.2026.3664118}}
\item
  \phantomsection\label{r25}{Z.-H. Zhou. A brief introduction to weakly
  supervised learning. \emph{National Science Review}, 5(1):44--53,
  2018. \href{https://doi.org/10.1093/nsr/nwx106}{10.1093/nsr/nwx106}}
\item
  \phantomsection\label{r26}{E. Romera, L. M. Bergasa, R. Arroyo. Need
  Data for Driver Behaviour Analysis? Presenting the Public
  UAH-DriveSet. \emph{IEEE ITSC}, 387--392, 2016.
  \href{https://doi.org/10.1109/ITSC.2016.7795584}{10.1109/ITSC.2016.7795584}}
\end{enumerate}
